\documentclass[11pt]{article} 
\usepackage{newpxtext,newpxmath}
\usepackage[dutch]{babel} 
\usepackage[a4paper, margin=2.5cm]{geometry} 
\usepackage{parskip} 
\author{}

\begin{document} 
\title{Samenwerkingsovereenkomst Co-auteurs (Zelfpublicatie)} 
\date{} 
\maketitle 
\section*{Partijen} 
\begin{itemize} 
    \item \textbf{[Lucas Arnoldus Hoogduin]}, wonende te 1185 MC Amstelveen, Burgemeester Rijnderslaan 332, hierna: \textit{Auteur 1} 
    \item \textbf{[Paul Touw]}, wonende te 8032 NL Zwolle, Deurzerdiep 29, hierna: \textit{Auteur 2} 
\end{itemize} 
Hierna gezamenlijk: \textit{Partijen} 

\section{Doel van de samenwerking} Partijen hebben gezamenlijk een boek geschreven getiteld: 
\begin{center} 
    \textbf{"Audit Data Analysis -- Volume 1"} 
\end{center} 
en exploiteren dit werk in eigen beheer (hierna: \textit{het Werk}), zowel fysiek als digitaal. 

\section{Auteursrechten} 

\begin{enumerate} 
    \item Het auteursrecht op het Werk berust gezamenlijk bij Partijen. 
    \item Beide auteurs hebben gelijke rechten, tenzij schriftelijk anders overeengekomen. 
    \item Besluiten over exploitatie worden gezamenlijk genomen. 
\end{enumerate} 

\section{Inkomsten en verdeling} 

\begin{enumerate} 
    \item Onder inkomsten vallen alle opbrengsten uit het Werk. 
    \item Netto-opbrengsten worden als volgt verdeeld: 
    \begin{itemize} 
        \item Auteur 1: \textbf{66 2/3 \%}
        \item Auteur 2: \textbf{33 1/3 \%}
    \end{itemize} 
\end{enumerate} 

\section{Kosten} 

\begin{enumerate} 
    \item Onder kosten vallen o.a. productie-, marketing- en distributiekosten. 
    \item Kosten worden gedragen in verhouding tot de inkomstenverdeling, tenzij anders afgesproken. 
\end{enumerate} 

\section{Beheer en management fee} 

\begin{enumerate} 
    \item Indien één partij verantwoordelijk is voor beheer (zoals website, administratie en marketing), ontvangt deze een management fee van \textbf{1 \%} van de bruto opbrengsten. 
    \item Deze fee wordt als kostenpost verwerkt vóór verdeling van de winst. 
\end{enumerate} 

\section{Jaarlijkse afrekening} 

\begin{enumerate} 
    \item Jaarlijks wordt per 31 december een overzicht opgesteld van opbrengsten en kosten. 
    \item Afrekening vindt plaats binnen 2 maanden na afloop van het jaar. 
    \item Beide Partijen hebben recht op inzage in de administratie. 
\end{enumerate} 

\section{Besluitvorming} Besluiten worden bij voorkeur unaniem genomen. Bij verschil van mening treden Partijen in overleg. 

\section{Duur en beëindiging} 
\begin{enumerate} 
    \item Deze overeenkomst wordt aangegaan voor onbepaalde tijd. 
    \item Opzegging geschiedt schriftelijk met inachtneming van de exitregeling. 
\end{enumerate} 

\section{Exit en uitkoop} 

\begin{enumerate} 
    \item Een partij kan de samenwerking schriftelijk opzeggen met een opzegtermijn van 3 maanden. 
    \item De andere partij heeft het recht het Werk voort te zetten en het aandeel over te nemen. 
    \item De uitkoopprijs wordt vastgesteld op:\\ \textbf{gemiddelde netto jaarwinst (laatste 2 jaar) × 2}, \\naar rato van het aandeel. 
    \item Betaling vindt plaats binnen 6 maanden, tenzij schriftelijk anders overeengekomen. 
    \item Bij geschil over de waarde wordt een onafhankelijke deskundige benoemd; diens oordeel is bindend. Kosten worden gelijk verdeeld. 
    \item De uittredende partij mag het Werk niet zelfstandig exploiteren, maar behoudt naamsvermelding als co-auteur. 
\end{enumerate} 

\section{Overdracht van rechten} Overdracht aan derden is alleen toegestaan met schriftelijke toestemming van beide Partijen. 

\section{Geschillen} 

\begin{enumerate} 
    \item Partijen proberen geschillen eerst onderling op te lossen. 
    \item Indien nodig wordt een mediator ingeschakeld. 
    \item Indien geen oplossing wordt bereikt, is de bevoegde rechter in Amsterdam bevoegd. 
\end{enumerate} 
        
\section{Slotbepalingen} 
    \begin{enumerate} 
        \item Wijzigingen zijn alleen geldig indien schriftelijk vastgelegd en door beide Partijen ondertekend. 
        \item Deze overeenkomst vervangt alle eerdere afspraken. 
\end{enumerate} 

\vspace{1cm} 

\textbf{Ondertekening} 

\vspace{0.5cm} Plaats: \rule{6cm}{0.4pt} Datum: \rule{6cm}{0.4pt} \vspace{1cm} 

\begin{tabular}{p{7cm} p{7cm}} \textbf{Auteur 1} & \textbf{Auteur 2} \\ 
    
\vspace{2cm} & \vspace{2cm} \\ 

Naam: Lucas Arnoldus Hoogduin & Naam: Paul Touw \\ 
Handtekening: & Handtekening:\\   
\vspace{2cm} & \vspace{2cm} \\ 
\rule{6cm}{0.4pt} & \rule{6cm}{0.4pt} \\ 

\end{tabular} 

\end{document}
